% mazzi: 34
\chapter{Il modello prestativo nel gioco del calcio}
\label{cap:modelloprestativo}

È attraverso la conoscenza del \textbf{modello prestativo} del calciatore che si riesce a progettare
e programmare l'intero percorso di allenamento. Il capitolo raccoglie ciò che la ricerca ha
stabilito sulla prestazione in gara --- fatica, risposta fisiologica, profilo di movimento, distanze
percorse --- e si chiude sul passaggio dal dato prestativo alla programmazione del carico.

% ---------------------------------------------------------------------------
\section{Il modello prestativo del calciatore}

\rubrica{Concetti da ricordare}
\begin{itemize}
  \item stagione agonistica con \textbf{50/60 partite ufficiali};
  \item ogni partita dura \textbf{90$'$ più i recuperi}, per un totale che arriva oltre i
    \textbf{100$'$} giocati;
  \item ogni partita è \textbf{unica} e quindi \textbf{non riproducibile in allenamento}.
\end{itemize}

\rubrica{Sforzo tipico del calciatore}
\begin{itemize}
  \item risponde a stimoli di tipo \textbf{intermittente};
  \item spesso \textbf{massimali};
  \item per lo più \textbf{brevi e ripetuti};
  \item \textbf{variabili} e anche \textbf{imprevedibili};
  \item di natura \textbf{alattacida e lattacida}.
\end{itemize}

\rubrica{Modello fisiologico}
Resistente alla capacità di \textbf{muoversi}, di \textbf{accelerare} e \textbf{decelerare} in spazi
e tempi brevissimi, con successioni imprevedibili e variabili nel tempo, intervallate da
\textbf{pause spesso minime} e comunque \textbf{sempre irregolari}.

\rubrica{Identità fisiologico-organica}
\begin{itemize}
  \item \textbf{discreta} capacità aerobica;
  \item \textbf{buona} capacità esplosiva;
  \item \textbf{ottima} resistenza allo sprint.
\end{itemize}

% ---------------------------------------------------------------------------
\section{La fatica in gara e la risposta fisiologica}

\subsection{Che tipo di fatica e da che cosa è causata}

\begin{itemize}
  \item è stabilito che \textbf{sprint}, \textbf{corsa ad alta intensità} e \textbf{distanza totale
    coperta} siano \textbf{più bassi nel 2\textsuperscript{o} tempo} rispetto al
    1\textsuperscript{o} (Reilly \& Thomas, 1979; Bangsbo, 1991--1994; Mohr, 2003; Carling \&
    Dupont, 2011);
  \item la corsa ad alta intensità è \textbf{fortemente condizionata dalle attività svolte durante
    il 1\textsuperscript{o} tempo};
  \item le \textbf{abilità tecniche non sono disturbate} dalla riduzione della capacità
    sforzo/lavoro;
  \item \textbf{salto}, \textbf{sprint} e \textbf{movimenti repentini} (COD) hanno invece una
    \textbf{diminuzione significativa}, valutati con test specifici prima e dopo la gara.
\end{itemize}

\subsection{La risposta cardiocircolatoria e metabolica}

\begin{itemize}
  \item il calcio è uno sport \textbf{intermittente}, dove il sistema aerobico è condizionato dalla
    \textbf{frequenza cardiaca}:
    \begin{itemize}
      \item in gara la FC si attesta attorno all'\textbf{85\% della massima};
      \item corrisponde a un consumo di ossigeno vicino al \textbf{70\% del
        $VO_2max$};
      \item la FC è \textbf{raramente al di sotto del 65\% della $FC_{max}$}: se ne deduce che la
        circolazione sanguigna per la cinematica muscolare sia continuamente più alta che a riposo,
        cioè che il \textbf{trasporto di ossigeno} sia elevato.
    \end{itemize}
  \item la \textbf{cinetica del consumo di ossigeno} durante le variazioni di movimento da bassa ad
    alta intensità dipende dalle \textbf{capacità ossidative della contrazione muscolare}
    (Bangsbo, 2010);
  \item in gara il calciatore esegue circa \textbf{150--250 azioni brevi e intense} $\rightarrow$
    il \textbf{tasso di scambio dell'energia anaerobica} è molto elevato, come pure l'utilizzo di
    \textbf{creatinfosfato};
  \item il \textbf{glicogeno muscolare} è il substrato più importante per la produzione di energia:
    la \textbf{fatica percepita} al termine della gara è correlata alla \textbf{deplezione di
    glicogeno} nelle fibre muscolari.
\end{itemize}

% ---------------------------------------------------------------------------
\section{Il profilo di movimento e le differenze fra ruoli}

L'analisi del profilo di movimento (Bloomfield, Polman \& O'Donoghue, 2007, sui giocatori della FA
Premier League) distingue i \textbf{movimenti significativi} (\textit{purposeful movements}, PM) per
tipo, intensità e direzione.

\rubrica{Che cosa emerge}
\begin{itemize}
  \item meno della \textbf{metà} dei movimenti significativi viene eseguita con \textbf{direzione
    frontale} (in avanti);
  \item l'attività a intensità \textbf{molto alta} è quasi il \textbf{doppio negli attaccanti
    rispetto ai difensori}, mentre fra attaccanti e centrocampisti i valori di intensità alta e
    molto alta sono \textbf{paragonabili};
  \item \textbf{difensori}:
    \begin{itemize}
      \item trascorrono un tempo significativamente \textbf{basso} in percentuale ai PM nella
        \textbf{corsa} e negli \textbf{sprint};
      \item di contro spendono \textbf{molto tempo saltando} rispetto agli altri ruoli;
      \item si muovono molto di più con \textbf{corse direzionali all'indietro}.
    \end{itemize}
  \item \textbf{centrocampisti}: effettuano \textbf{molti meno cambi di direzione} rispetto ad
    attaccanti e difensori.
\end{itemize}

\rubrica{La frequenza dei cambi di direzione}
La stragrande maggioranza dei cambi di direzione, come numero, avviene con angoli fra
\textbf{0--90$^\circ$} e fra \textbf{90--180$^\circ$}, sia a destra sia a sinistra, e questo vale
per tutti e tre i ruoli.

\begin{table}[htbp]
  \centering
  \small
  \renewcommand{\arraystretch}{1.3}
  \begin{tabularx}{\textwidth}{|l|X|X|X|X|c|}
    \hline
    \textbf{Variabile} & \textbf{Attaccanti} \newline ($n=19$) & \textbf{Centrocamp.} \newline
      ($n=18$) & \textbf{Difensori} \newline ($n=18$) & \textbf{Tutti} \newline ($n=55$) &
      \textbf{$p$} \\
    \hline
    0--90$^\circ$ destra    & 323,7 (105,1) & 248,3 (97,3) & 344,3 (91,0) & 305,8 (104,7) & ,010 \\
    \hline
    0--90$^\circ$ sinistra  & 302,2 (81,2) & 243,0 (93,5) & 364,3 (88,4) & 303,2 (99,3) & ,001 \\
    \hline
    90--180$^\circ$ destra  & 43,3 (15,6) & 49,3 (25,0) & 43,0 (16,8) & 45,2 (19,4) & ,898 \\
    \hline
    90--180$^\circ$ sinistra & 51,5 (13,9) & 47,0 (24,5) & 49,3 (21,4) & 49,3 (20,1) & ,578 \\
    \hline
    180--270$^\circ$ destra & 2,5 (4,2) & 4,7 (3,9) & 2,3 (3,0) & 3,2 (3,8) & ,098 \\
    \hline
    180--270$^\circ$ sinistra & 2,2 (3,6) & 3,0 (4,7) & 2,0 (2,9) & 2,4 (3,8) & ,926 \\
    \hline
    270--360$^\circ$ destra & 1,3 (2,5) & 0,7 (1,9) & 0,0 (0,0) & 0,7 (1,9) & ,126 \\
    \hline
    270--360$^\circ$ sinistra & 0,6 (1,9) & 2,3 (3,6) & 0,0 (0,0) & 1,0 (2,5) & ,015 \\
    \hline
    \textit{Swerve} destra  & 8,5 (8,3) & 5,7 (7,3) & 7,7 (6,4) & 7,3 (7,4) & ,424 \\
    \hline
    \textit{Swerve} sinistra & 12,0 (9,6) & 4,0 (6,5) & 9,3 (10,3) & 8,5 (9,4) & ,015 \\
    \hline
    \textbf{Totale}         & 748 (173) & 608 (207) & 822 (175) & 727 (203) & ,010 \\
    \hline
  \end{tabularx}
\end{table}

\rubrica{La conseguenza metodologica}
Queste differenze indicano che i calciatori devono essere allenati attraverso \textbf{programmi
specifici}, tenendo conto:
\begin{enumerate}
  \item del \textbf{ruolo} ricoperto in campo;
  \item delle \textbf{qualità fisiche} richieste;
  \item e soprattutto della \textbf{struttura fisica individuale} del calciatore.
\end{enumerate}

% ---------------------------------------------------------------------------
\section{Il profilo fisiologico e il quadro prestativo del gioco}

\subsection{Il lattato}

\begin{itemize}
  \item ``Malgrado l'alto numero di sprint il LA va molto di rado oltre i \textbf{6--7 mmol/L}''
    (Mombaerts, 1991);
  \item ``una delle conclusioni è di \textbf{evitare gli sforzi prolungati} atti a provocare una
    alta produzione di LA. Per questo si consiglia di \textbf{non proporre in allenamento i lavori
    ad alta intensità e di tempo $> 30''$} tendenti a stimolare la tolleranza ad alte
    concentrazioni di LA'' (Chatard, 1992).
\end{itemize}

\begin{table}[htbp]
  \centering
  \small
  \renewcommand{\arraystretch}{1.3}
  \begin{tabularx}{\textwidth}{|X|X|c|c|}
    \hline
    \textbf{Studio} & \textbf{Giocatori} & \textbf{Fine 1\textsuperscript{o} tempo} &
      \textbf{Fine 2\textsuperscript{o} tempo} \\
    \hline
    Agnevik, 1970 & prima divisione & --- & 10 \\
    \hline
    Ekblom, 1986 & prima divisione & 9,5 & 7,2 \\
    \hline
     & seconda divisione & 8 & 6,6 \\
    \hline
     & terza divisione & 5,5 & 4,2 \\
    \hline
     & quarta divisione & 4 & 3,9 \\
    \hline
    Mombaerts, 1991 & --- & --- & 6--7 \\
    \hline
    Bangsbo, 1993 & campionato danese & 2,6 & 2,7 \\
    \hline
    Sassi--Candel, 1998 & --- & --- & 4,2 \\
    \hline
    Nanni e coll., 2001 & Serie A & --- & 6,2 \\
    \hline
  \end{tabularx}
  \caption*{Concentrazioni medie di LA (mM) nel calcio.}
\end{table}

\rubrica{Lattato e tempi di recupero (Balsom)}
Il confronto fra esercitazioni a pari volume ma con recuperi diversi mostra quanto il
\textbf{recupero} governi la produzione di lattato:
\begin{itemize}
  \item \textbf{15 $\times$ 40 m con 30$''$ di recupero}: al termine si arriva attorno alle
    \textbf{16 mmol/L}, e nei minuti successivi la concentrazione \textbf{sale ancora}, fino a
    18--19 mmol/L $\rightarrow$ non è un andamento ottimale;
  \item \textbf{40 $\times$ 15 m con 30$''$ di recupero}: si resta attorno alle \textbf{6 mmol/L}
    al termine, e dopo 5 minuti il lattato è già stato smaltito;
  \item la differenza non sta quindi nel volume complessivo, ma nella \textbf{lunghezza della
    singola ripetizione} rispetto al recupero concesso.
\end{itemize}

\subsection{Gli sprint}

Il numero medio di sprint prodotti in gara in funzione della distanza (Dufour, 1990) descrive una
curva con un \textbf{massimo netto sui 15 m} circa, e un crollo rapido oltre i 25 m.

\begin{itemize}
  \item la distanza che conta per la prestazione calcistica va dai \textbf{5} ai
    \textbf{15--20--25 m}, non oltre;
  \item è l'indicazione operativa su quali distanze far percorrere in modalità sprint in
    allenamento.
\end{itemize}

\subsection{Accelerazioni e decelerazioni}

\begin{itemize}
  \item \textbf{accelerazioni}: durante la gara il \textbf{20\%} si sviluppa \textbf{oltre il 50\%
    della massima possibile}, a qualsiasi velocità; tutte le altre sono \textbf{accelerazioni
    moderate};
  \item \textbf{decelerazioni}: solo per il \textbf{5\%} il calciatore frena in modo
    \textbf{massimale}; molte di più sono le \textbf{decelerazioni medie};
  \item a parità di velocità le decelerazioni avvengono \textbf{molto più rapidamente} delle
    accelerazioni:
    \begin{itemize}
      \item frenare da 20 km/h può avvenire in \textbf{500 ms};
      \item accelerare fino a 20 km/h richiede circa \textbf{1,5--2$''$}.
    \end{itemize}
\end{itemize}

\subsection{La forza applicata nelle diverse tipologie di spostamento}

\begin{itemize}
  \item \textbf{accelerazione}: forza \textbf{elevata} e tempo di applicazione \textbf{prolungato};
  \item \textbf{corsa a soglia}: forza \textbf{più bassa};
  \item \textbf{sprint}: forza applicata \textbf{più alta} e nel \textbf{tempo più breve};
  \item \textbf{corsa a ritmo elevato e prolungata}: la forza \textbf{diminuisce}, compensata da
    una \textbf{frequenza dei passi più alta}.
\end{itemize}

\rubrica{Le espressioni di forza nei 100 m (Carlo Vittori)}
La combinazione delle diverse espressioni di forza lungo una gara di 100 m distingue due fasi:
\begin{itemize}
  \item \textbf{fase di accelerazione}, con la \textbf{messa in moto} nei primi 10 m: domina la
    \textbf{forza esplosiva} (forza attiva) e la \textbf{forza elastica};
  \item \textbf{fase lanciata}: domina la \textbf{forza reattiva}, cioè la \textit{stiffness}
    eccentrico-riflessa, che poggia sulla \textbf{capacità contrattile};
  \item la \textbf{componente di forza contrattile} attraversa tutta la gara e per il calciatore è
    quella sviluppata fra i \textbf{10 e i 35 m};
  \item ciò che rileva di più per il calciatore sono comunque i \textbf{primi 15 m} della
    prestazione di sprint.
\end{itemize}

\subsection{I cambi di direzione}

\begin{itemize}
  \item sono presenti in modo elevato: circa \textbf{1000} con angoli superiori ai
    \textbf{30$^\circ$}, e oltre \textbf{800} sviluppati anche a potenze superiori a
    \textbf{20 W/kg};
  \item in pratica il calciatore sviluppa un cambio di direzione $> 30^\circ$ e a $> 20$ W/kg ogni
    \textbf{17--18$''$};
  \item gli \textbf{studi futuri} dovranno tendere a mettere in rapporto i cambi di direzione sia
    con la \textbf{potenza} sia con la \textbf{velocità}: serve però un progresso tecnologico dei
    \textbf{GPS}, che attualmente non riescono ad analizzare spostamenti così repentini e con
    direzioni variabili.
\end{itemize}

% ---------------------------------------------------------------------------
\section{Il modello di prestazione secondo Di Salvo}

Studio su \textbf{563 giocatori} della \textbf{Premier League} inglese, stagioni
\textbf{2003--2006}.

\rubrica{Le sigle}
\begin{itemize}
  \item \textbf{AIT}: alta intensità totale;
  \item \textbf{AITP}: alta intensità totale \textbf{con} possesso palla;
  \item \textbf{AITSP}: alta intensità totale \textbf{senza} possesso palla.
\end{itemize}

\rubrica{I risultati}
\begin{itemize}
  \item l'\textbf{attività di sprint} (velocità $> 25,2$ km/h) dei \textbf{centrocampisti
    laterali}, degli \textbf{attaccanti} e dei \textbf{difensori laterali} è significativamente
    maggiore rispetto a quella dei \textbf{centrocampisti centrali} e dei \textbf{difensori
    centrali};
  \item l'\textbf{attività totale ad alta intensità} ($> 19,8$ km/h) effettuata dagli
    \textbf{attaccanti con la propria squadra in possesso} è maggiore rispetto a quella di tutti
    gli altri ruoli;
  \item di contro, la stessa attività \textbf{con la propria squadra non in possesso} mostra gli
    attaccanti come i \textbf{meno attivi} rispetto agli altri ruoli.
\end{itemize}

% ---------------------------------------------------------------------------
\section{L'analisi prestativa con la match analysis}

L'impostazione è quella di \textbf{Daniele Tognaccini}, responsabile del \textit{Milan Lab}.

\begin{itemize}
  \item gli \textbf{sport di situazione} si caratterizzano per la \textbf{difficoltà di codificare}
    in modo standardizzato e costante la prestazione;
  \item fra gli strumenti più utilizzati a livello di élite per valutare l'impegno dell'atleta in
    gara e in allenamento c'è il sistema di \textbf{\textit{time motion analysis}}, che permette di
    monitorare la prestazione sulla base di alcuni parametri;
  \item resta ancora \textbf{non definito il costo energetico} determinato dai \textbf{cambi di moto
    della massa del corpo}:
    \begin{itemize}
      \item salti verticali;
      \item cambi di direzione;
      \item cambi di velocità in accelerazione e decelerazione;
      \item contatti con l'avversario;
      \item tutti i movimenti compiuti in lateralità.
    \end{itemize}
\end{itemize}

\rubrica{La conseguenza: l'individualizzazione}
\begin{itemize}
  \item risulterebbe \textbf{difficile}, al termine di una gara, determinare se il giocatore
    analizzato \textbf{è in forma o no} senza tener conto della \textbf{specificità della
    situazione};
  \item ciascuna analisi, per avere un valore maggiormente significativo, dovrà tener conto:
    \begin{itemize}
      \item del \textbf{singolo atleta} come soggetto specifico, \textbf{unico e irripetibile};
      \item di ciascuna \textbf{squadra} come \textbf{sommatoria di soggetti unici e irripetibili}.
    \end{itemize}
  \item è quindi essenziale un'\textbf{osservazione specifica} della squadra e dei singoli, con
    riferimento a un modello di prestazione ``ideale'' \textbf{assolutamente individualizzato}.
\end{itemize}

% ---------------------------------------------------------------------------
\section{La distanza totale percorsa}

\subsection{Che cosa si misura e con che cosa}

\begin{itemize}
  \item l'\textbf{intensità della corsa} in gara è indicata dalla distanza percorsa attraverso
    \textbf{differenti zone di velocità}, dalla \textbf{velocità media} e dalla \textbf{distanza
    totale percorsa} (TDC) in partita;
  \item gli strumenti impiegati sono diversi:
    \begin{itemize}
      \item \textit{video-based time-motion analysis};
      \item \textit{semi-automatic multiple-camera system};
      \item GPS.
    \end{itemize}
  \item poiché strumentazioni diverse producono \textbf{larghe differenze} --- per esempio proprio
    sulla TDC --- nel comparare i dati ottenuti va usata \textbf{molta cautela}
    nell'interpretazione.
\end{itemize}

\begin{table}[htbp]
  \centering
  \small
  \renewcommand{\arraystretch}{1.3}
  \begin{tabularx}{\textwidth}{|l|X|c|X|}
    \hline
    \textbf{Autore} & \textbf{Livello} & \textbf{TDC (km)} & \textbf{Metodo} \\
    \hline
    Bangsbo (1994) & internazionali danesi & 10,55 & \textit{time motion analysis} \\
    \hline
    Helgerud et al.\ (2001) & internazionali junior norvegesi & 9,1 & \textit{time motion analysis} \\
    \hline
    Mohr et al.\ (2004) & professionisti danesi & 10,33 & \textit{time motion analysis} \\
    \hline
    Di Salvo et al.\ (2007) & professionisti spagnoli & 11,39 & \textit{video tracking} \\
    \hline
    Barros et al.\ (2007) & professionisti brasiliani & 10,01 & \textit{video tracking} \\
    \hline
    Dellal et al.\ (2010, 2011) & prima divisione (Francia, UK, Spagna, Italia) & 10,42--11,78 &
      \textit{video tracking} \\
    \hline
  \end{tabularx}
  \caption*{Distanza totale media percorsa da un giocatore in una gara ufficiale.}
\end{table}

\subsection{La ripartizione per intensità}

\rubrica{Il quadro d'insieme}
\begin{itemize}
  \item \textbf{84,8\%} del tempo camminando e correndo lentamente;
  \item \textbf{11,3\%} del tempo ad attività intensa (cambi di passo e sprint);
  \item rapporto di \textbf{1:7} fra alta e bassa attività di corsa;
  \item l'\textbf{88\%} delle gare si basa su componenti di \textbf{resistenza} e il rimanente
    \textbf{12\%} su componenti \textbf{intense};
  \item è però proprio quell'11--12\% il momento in cui si estrinseca la \textbf{massima
    performance} del calciatore, non l'85\% trascorso camminando o correndo lentamente.
\end{itemize}

\rubrica{La durata delle azioni e dei recuperi (Orenduff et al., 2010)}
\begin{itemize}
  \item \textbf{durata dei movimenti}: 43\% meno di 6$''$; 23\% fra 6 e 9$''$; 13\% fra 9 e
    12$''$; 9\% fra 12 e 15$''$;
  \item \textbf{durata dei recuperi}: 53\% meno di 6$''$; 22\% fra 6 e 9$''$; 9\% fra 9 e 12$''$;
    5\% fra 12 e 15$''$.
\end{itemize}

\rubrica{Il numero delle azioni}
\begin{itemize}
  \item i giocatori compiono fra \textbf{1000 e 1400 azioni di breve durata} in una gara, della
    durata di \textbf{2--4$''$} (Stolen et al., 2005), di cui circa \textbf{220 ad alta intensità}
    (Mohr, 2003);
  \item il giocatore compie un'\textbf{azione diversa ogni 4--6$''$} per tutta la durata della gara
    (Bangsbo, 1994).
\end{itemize}

\subsection{Le differenze fra ruoli}

\begin{table}[htbp]
  \centering
  \small
  \renewcommand{\arraystretch}{1.3}
  \begin{tabularx}{\textwidth}{|X|c|c|c|}
    \hline
    \textbf{Azione} & \textbf{Difensori (m)} & \textbf{Centrocampisti (m)} &
      \textbf{Attaccanti (m)} \\
    \hline
    Camminata & 2691 & 2451 & 3007 \\
    \hline
    Camminata all'indietro & 562 & 569 & 526 \\
    \hline
    Corsa lenta (\textit{jog}) & 3869 & 4892 & 2605 \\
    \hline
    Corsa lenta all'indietro & 276 & 300 & 68 \\
    \hline
    Corsa lenta laterale & 362 & 319 & 73 \\
    \hline
    \textit{Cruise} & 701 & 1110 & 900 \\
    \hline
    Sprint & 231 & 316 & 557 \\
    \hline
  \end{tabularx}
  \caption*{Distanza media percorsa per ruolo dai nazionali sudamericani (Rienzi et al., 2000).}
\end{table}

Il dato che emerge è che gli \textbf{attaccanti} percorrono in sprint distanze molto più elevate
degli altri ruoli: \textbf{557 m} contro i \textbf{231 m} dei difensori, cioè meno della metà.

\begin{table}[htbp]
  \centering
  \small
  \renewcommand{\arraystretch}{1.3}
  \begin{tabularx}{\textwidth}{|X|c|c|c|c|}
    \hline
    & \textbf{Dif.\ centrali} & \textbf{Dif.\ laterali} & \textbf{Centrocamp.} &
      \textbf{Attaccanti} \\
    \hline
    TDC (m) & 9995 & 11232 & 11748 & 10233 \\
    \hline
    TDC camminando (0,7--2 km/h) & 3846 & 3504 & 3341 & 3844 \\
    \hline
    TDC in corsa lenta (7,2--14,4 km/h) & 1458 & 1601 & 1726 & 1361 \\
    \hline
    TDC in \textit{cruising} (14,4--19,8 km/h) & 278 & 211 & 467 & 321 \\
    \hline
    TDC ad alta intensità (19,8--25,2 km/h) & 76 & 123 & 118 & 95 \\
    \hline
    Numero di sprint ($> 25,2$ km/h) & 18 & 31 & 24 & 27 \\
    \hline
  \end{tabularx}
  \caption*{Distanza media percorsa per ruolo dai professionisti inglesi (Rampinini et al., 2007).}
\end{table}

\rubrica{Premier League e Liga a confronto (Dellal, 2011)}
Il confronto fra i due campionati, ruolo per ruolo, conferma due cose:
\begin{itemize}
  \item la \textbf{TDC} è ormai stabilmente sopra gli 11 km, e nella Liga i \textbf{difensori
    centrali} arrivano a circa \textbf{14.096 m}, contro i \textbf{10.064 m} dei difensori
    laterali e i \textbf{10.718 m} degli attaccanti centrali;
  \item la \textbf{distanza ad alta intensità in fase di non possesso} conferma il quadro di Di
    Salvo: gli \textbf{attaccanti centrali} ne percorrono pochissima --- circa \textbf{93 m} nella
    Liga e \textbf{101 m} in Premier League --- a fronte di valori molto più alti nei ruoli
    arretrati.
\end{itemize}

\subsection{Primo e secondo tempo}

\begin{table}[htbp]
  \centering
  \small
  \renewcommand{\arraystretch}{1.3}
  \begin{tabularx}{\textwidth}{|X|X|c|c|c|c|}
    \hline
    \textbf{Autore} & \textbf{Livello} & \textbf{1\textsuperscript{o} t.} &
      \textbf{2\textsuperscript{o} t.} & \textbf{TDC} & \textbf{Diff.} \\
    \hline
    Rienzi et al.\ (2000) & internazionali sudamericani & 4,60 & 4,41 & 8,63 & $-4\%$ \\
    \hline
    Mohr et al.\ (2003) & professionisti danesi & 5,51 & 5,35 & 10,86 & $-3\%$ \\
    \hline
    Di Salvo et al.\ (2007) & professionisti spagnoli & 5,70 & 5,68 & 11,39 & $-1\%$ \\
    \hline
    Rampinini et al.\ (2007) & professionisti inglesi & --- & --- & 10,86 & $-8\%$ \\
    \hline
    Barros et al.\ (2007) & professionisti brasiliani & 5,17 & 4,80 & 10,01 & $-7\%$ \\
    \hline
  \end{tabularx}
  \caption*{Variazioni della distanza totale percorsa (km) fra primo e secondo tempo.}
\end{table}

Il decremento fra i due tempi è la misura più immediata della \textbf{fatica}, e varia molto con il
livello e il campionato: minimo nei professionisti spagnoli ($-1\%$), massimo nei brasiliani e negli
inglesi ($-7$/$-8\%$).

\begin{table}[htbp]
  \centering
  \small
  \renewcommand{\arraystretch}{1.3}
  \begin{tabularx}{\textwidth}{|X|c|c|}
    \hline
    & \textbf{TDC 1\textsuperscript{o} tempo (m)} & \textbf{TDC 2\textsuperscript{o} tempo (m)} \\
    \hline
    \textbf{Totale} & 5709 & 5684 \\
    \hline
    0--11 km/h & 3496 & 3535 \\
    \hline
    11,1--14 km/h & 851 & 803 \\
    \hline
    14,1--19 km/h & 894 & 865 \\
    \hline
    19,1--23 km/h & 304 & 301 \\
    \hline
    $> 23$ km/h & 165 & 172 \\
    \hline
    Con la palla & 104 & 109 \\
    \hline
  \end{tabularx}
  \caption*{Distanza media percorsa nei due tempi alle diverse velocità (Di Salvo et al., 2007).}
\end{table}

La lettura per zone corregge il quadro: alle andature basse la distanza è \textbf{sostanzialmente
identica} nei due tempi (circa 3500 m fra 0 e 11 km/h), il calo si concentra nelle fasce
intermedie, e \textbf{sopra i 23 km/h il secondo tempo è addirittura superiore} al primo.

% ---------------------------------------------------------------------------
\section{Dal modello prestativo al carico di allenamento}

\rubrica{Il principio}
La conoscenza dettagliata dello \textbf{stress fisiologico} --- \textbf{carico esterno} e
\textbf{carico interno} --- prodotto da ogni singola esercitazione consente una loro
\textbf{corretta collocazione} nel corso delle unità di allenamento.

\subsection{Il confronto con le altre squadre}

Un esempio applicativo viene dal monitoraggio di una squadra di prima divisione rumena, con i dati
\textit{Instat Fitness Data} del girone di andata 2014, che confrontano ogni squadra con la media
del campionato su distanza totale e sulle singole fasce di velocità:

\begin{itemize}
  \item \textbf{camminata}: 0--2 m/s (0--7,2 km/h);
  \item \textbf{corsa leggera}: 2--4 m/s (7,2--14,4 km/h);
  \item \textbf{corsa quantitativa}: 4--5,5 m/s (14,4--19,8 km/h);
  \item \textbf{accelerazioni}: 5,5--7 m/s (19,8--25,2 km/h);
  \item \textbf{sprint}: $> 7$ m/s ($> 25,2$ km/h).
\end{itemize}

\rubrica{Che cosa se ne ricava}
\begin{itemize}
  \item sulla \textbf{distanza totale} del primo tempo la squadra era in linea con la media del
    campionato (57.256 m contro 57.286 m, cioè appena 30 m in meno);
  \item sulle fasce di \textbf{bassa e media intensità} i valori erano a favore della squadra;
  \item sulla \textbf{distanza in sprint} ($> 25,2$ km/h) i valori erano invece, seppure di poco,
    \textbf{al di sotto della media} delle altre squadre;
  \item lettura: la squadra aveva \textbf{buone capacità di resistenza}, estrinsecabili per tutto
    l'arco della partita, ma \textbf{non riusciva a produrre picchi prestativi} nelle attività di
    sprint $\rightarrow$ dato fondamentale per la progettazione e la programmazione del
    \textbf{girone di ritorno}.
\end{itemize}

\subsection{L'indice di usura (\texorpdfstring{\textit{stress index}}{stress index})}

\textbf{Stress index} (indice di usura): rapporto fra i minuti effettivamente giocati e i minuti
totali di partita disputati dalla squadra.

\[
  \textit{Stress index} = \frac{\textrm{minuti giocati}}{\textrm{minuti partite}} \times 100
\]

\begin{itemize}
  \item è semplicissimo da calcolare e permette al preparatore di conoscere, \textbf{partita dopo
    partita}, il \textbf{minutaggio} e quindi lo stress accumulato da ciascun giocatore;
  \item sul grafico dell'andamento individuale si tracciano due soglie di riferimento:
    \begin{itemize}
      \item \textbf{75\%} (linea rossa);
      \item \textbf{50\%} (linea verde).
    \end{itemize}
  \item la lettura si traduce direttamente in due programmazioni opposte:
    \begin{itemize}
      \item giocatori \textbf{sopra il 75\%}: per loro si deve pensare piuttosto al
        \textbf{recupero} e al \textbf{ripristino delle energie} spese nella prestazione agonistica;
      \item giocatori \textbf{sotto il 50\%} --- utilizzati poco o per nulla: serve un programma di
        allenamento che renda la loro prestazione \textbf{idonea} nel momento in cui l'allenatore
        li chiama in campo.
    \end{itemize}
  \item è quindi uno strumento per la progettazione dei carichi \textbf{di squadra} ma soprattutto
    dei carichi \textbf{individuali}.
\end{itemize}
